\documentclass{article}
\usepackage[brazil]{babel}
\usepackage{booktabs}
\usepackage{makecell}
\usepackage{ragged2e}

\begin{document}
	
	\begin{table}[htbp]
		\centering
		\caption{Trabalhos analisados na revisão}
		\label{tab:trabalhos}
		\footnotesize
		\setlength{\tabcolsep}{4pt}
		\renewcommand{\arraystretch}{1.2}
		\begin{tabular}{p{7cm} c c c}
			\toprule
			\textbf{Título} & \textbf{Ano} & \textbf{GitHub} & \textbf{LMS} \\
			\midrule
			\makecell[l]{Identificação de Pesquisas e Análise de \\ Algoritmos de Clusterização para a \\ Descoberta de Perfis de Engajamento} & 2022 & Não & Sim \\
			\midrule
			\makecell[l]{A Educação no contexto da pandemia \\ de COVID-19: uma revisão sistemática \\ de literatura} & 2020 & Não & Não \\
			\midrule
			\makecell[l]{Integração de Tecnologias Digitais à \\ Linguagem Oral e Escrita: um estudo \\ de caso no Ensino Fundamental I} & 2018 & Não & Não \\
			\midrule
			\makecell[l]{Sobre a necessidade de recursos \\ educacionais para o ensino do \\ Pensamento Computacional na \\ Educação Básica Brasileira: discussão \\ e concepção de Repositório \\ Educacional do Pensamento \\ Computacional} & 2021 & Não & Não \\
			\midrule
			\makecell[l]{Small Open Online Course com \\ Professores do Ensino Médio: desafios \\ para integrar REA nos materiais e \\ atividades didáticas} & 2016 & Não & Não \\
			\midrule
			\makecell[l]{Impacto da descontinuidade da \\ tecnologia Flash na disponibilidade \\ dos Objetos de Aprendizagem em \\ repositórios educacionais} & 2021 & Não & Não \\
			\midrule
			\makecell[l]{Versionamento de Projeto na Prática \\ com Git e GitHub: Um Relato de \\ Experiência do Curso Ofertado pelo \\ Projeto LearningLab No Interior \\ Cearense} & 2024 & Sim & Não \\
			\midrule
			\makecell[l]{Jogos Digitais sob a perspectiva da \\ Engenharia de Software: um \\ mapeamento sistemático da literatura} & 2024 & Não & Não \\
			\midrule
			\makecell[l]{Implementação das Leis 10.639/03 e \\ 11.645/08 no Ensino de Computação: \\ Um Mapeamento da Literatura em \\ Bases Nacionais} & 2025 & Não & Não \\
			\midrule
			\makecell[l]{OntoObADi: Uma Ontologia para \\ Repositório Web Semântico de \\ Objetos de Aprendizagem Digitais} & 2025 & Não & Não \\
			\midrule
			\makecell[l]{Uma API para Gerenciamento de \\ Objetos de Aprendizagem Digitais \\ Baseada em Web Semântica} & 2025 & Não & Não \\
			\midrule
			\makecell[l]{Commit Explorer: Uma Solução Open \\ Source para Extração e Avaliação de \\ Commits e Códigos} & 2025 & Sim & Não \\
			\midrule
			\makecell[l]{ObserveUnB - Um portal Web de rede \\ social científica} & 2012 & Não & Não \\
			\midrule
			\makecell[l]{Práticas Educacionais Abertas \\ Mediadas por Tecnologia na Formação \\ Inicial Docente: em Direção à \\ Disseminação e Efetivação da \\ Cultura Aberta na Educação Básica} & 2025 & Não & Não \\
			\midrule
			\makecell[l]{Uma Análise Comparativa entre \\ Repositórios de Recursos \\ Educacionais Abertos para a \\ Educação Básica} & 2021 & Não & Não \\
			\midrule
			\makecell[l]{Usando heurísticas de Nielsen para \\ avaliar objetos de aprendizagem e \\ softwares educacionais: um estudo \\ Exploratório na área de Matemática \\ para ensino superior} & 2012 & -- & -- \\
			\bottomrule
		\end{tabular}
	\end{table}
	
\end{document}